\section{Proofs of the main theoretical results}\label{supp:proofs}

\subsection{Smoothing target approximation}\label{supp:target-proof}

\begin{theorem}[Target approximation]\label{thm:supp-target}
Under Assumption~\ref{ass:supp-target},
\[
 \norm{\beta_h^{\star}-\beta^{\star}}_{\infty}\leq C_HC h^2.
\]
If the difference is supported on a set of order $s$, then
$\norm{\beta_h^{\star}-\beta^{\star}}_2=O(\sqrt{s}\,h^2)$.
\end{theorem}

\begin{proof}
Let $d_h=\beta_h^{\star}-\beta^{\star}$. By the first-order condition and the integral mean-value identity,
\[
 0=\nabla R_h(\beta_h^{\star})
 =\nabla R_h(\beta^{\star})+\overline H_hd_h.
\]
Thus $d_h=-\overline H_h^{-1}\nabla R_h(\beta^{\star})$, and
\[
 \norm{d_h}_{\infty}
 \leq\norm{\overline H_h^{-1}}_{L_1}
 \norm{\nabla R_h(\beta^{\star})}_{\infty}
 \leq C_HC h^2.
\]
The $\ell_2$ statement follows from
$\norm{d_h}_2\leq\sqrt{|\operatorname{supp}(d_h)|}\norm{d_h}_{\infty}$.
\end{proof}

\subsection{Penalised profile rate}\label{supp:estimation-proof}

\begin{theorem}[Penalised estimation]\label{thm:supp-estimation}
Under Assumption~\ref{ass:supp-estimation}, choose
$\lambda_{\beta}\geq2\norm{g_{n,h}(\beta_h^{\star})}_{\infty}$. Then
\[
 \norm{\widehat\beta-\beta_h^{\star}}_2
 \leq C\sqrt{s}\lambda_{\beta},
 \qquad
 \norm{\widehat\beta-\beta_h^{\star}}_1
 \leq Cs\lambda_{\beta}.
\]
The same orders hold for a finite block of unpenalised covariates after augmenting the restricted-curvature condition by that block.
\end{theorem}

\begin{proof}
Set $\Delta=\widehat\beta-\beta_h^{\star}$ and $S=S_h$. The basic inequality is
\begin{equation}\label{eq:supp-basic}
 L_{n,h}^{\mathrm{prof}}(\beta_h^{\star}+\Delta)
 -L_{n,h}^{\mathrm{prof}}(\beta_h^{\star})
 +\lambda_{\beta}
 \{\norm{\beta_h^{\star}+\Delta}_1-\norm{\beta_h^{\star}}_1\}\leq0.
\end{equation}
Restricted curvature and the score bound give
\[
 L_{n,h}^{\mathrm{prof}}(\beta_h^{\star}+\Delta)
 -L_{n,h}^{\mathrm{prof}}(\beta_h^{\star})
 \geq g_{n,h}(\beta_h^{\star})^{\mathsf T}\Delta
 +\frac{\kappa}{2}\norm{\Delta}_2^2,
\]
and
$|g_{n,h}(\beta_h^{\star})^{\mathsf T}\Delta|
\leq(\lambda_{\beta}/2)\norm{\Delta}_1$.
The decomposability inequality
\[
 \norm{\beta_h^{\star}+\Delta}_1-\norm{\beta_h^{\star}}_1
 \geq\norm{\Delta_{S^c}}_1-\norm{\Delta_S}_1
\]
inserted into \eqref{eq:supp-basic} yields
\[
 \frac{\kappa}{2}\norm{\Delta}_2^2
 +\frac{\lambda_{\beta}}{2}\norm{\Delta_{S^c}}_1
 \leq\frac{3\lambda_{\beta}}{2}\norm{\Delta_S}_1.
\]
Hence $\Delta\in\mathcal C(S,3)$ and
$\kappa\norm{\Delta}_2^2/2
\leq3\lambda_{\beta}\sqrt{s}\norm{\Delta}_2/2$.
This proves the $\ell_2$ bound. The cone inequality then gives
$\norm{\Delta}_1\leq4\norm{\Delta_S}_1
\leq4\sqrt{s}\norm{\Delta}_2$, proving the $\ell_1$ bound.
\end{proof}

\subsection{Row-wise Dantzig inversion}\label{supp:precision-proof}

\begin{theorem}[Precision-row rate]\label{thm:supp-precision}
Under Assumption~\ref{ass:supp-precision}, any solution of the row-wise Dantzig program satisfies
\[
 \norm{\widehat\omega_k-\omega_{k,h}}_{\infty}=O_p(\mu_n),
 \qquad
 \norm{\widehat\omega_k-\omega_{k,h}}_1
 =O_p(s_{\omega,k}\mu_n).
\]
\end{theorem}

\begin{proof}
Write $H=H_h$, $\widehat H=\widehat H_{\mathrm{eff}}(\widehat\beta)$,
$\omega=\omega_{k,h}$, and $u=\widehat\omega_k-\omega$.
Because $H\omega=e_k$,
\[
 \norm{\widehat H\omega-e_k}_{\infty}
 \leq\norm{\widehat H-H}_{\max}\norm{\omega}_1
 \leq C_{\omega}\delta_{n,h}\leq\mu_n.
\]
Thus $\omega$ is feasible. The $\ell_1$ objective implies
$\norm{u_{T^c}}_1\leq\norm{u_T}_1$ for $T=\operatorname{supp}(\omega)$, and therefore
$\norm{u}_1\leq2s_{\omega,k}\norm{u}_{\infty}$.
Since both $\widehat\omega_k$ and $\omega$ are feasible,
$\norm{\widehat Hu}_{\infty}\leq2\mu_n$. Hence
\begin{align*}
 \norm{u}_{\infty}
 &\leq\norm{H^{-1}}_{L_1}\norm{Hu}_{\infty}\\
 &\leq C_{\Omega}
 \{\norm{(H-\widehat H)u}_{\infty}+\norm{\widehat Hu}_{\infty}\}\\
 &\leq C_{\Omega}\{2s_{\omega,k}\delta_{n,h}\norm{u}_{\infty}+2\mu_n\}.
\end{align*}
The first term is absorbed because $s_{\omega,k}\delta_{n,h}=o(1)$.
This proves the $\ell_{\infty}$ rate, and the cone bound gives the $\ell_1$ rate.
\end{proof}

\subsection{Cluster Bahadur representation}\label{supp:bahadur-proof}

\begin{theorem}[Bahadur representation]\label{thm:supp-bahadur}
Under Assumptions~\ref{ass:supp-estimation}--\ref{ass:supp-expansion},
\[
 \sqrt n(\widetilde\beta_k-\beta_{h,k}^{\star})
 =\frac1{\sqrt n}\sum_{i=1}^n
 \omega_{k,h}^{\mathsf T}S_{i,h}(\beta_h^{\star})+o_p(1).
\]
\end{theorem}

\begin{proof}
Let $g^{\star}=g_{n,h}(\beta_h^{\star})$ and
$\Delta=\widehat\beta-\beta_h^{\star}$. By
\eqref{eq:supp-score-expansion},
\begin{align*}
 \widetilde\beta_k-\beta_{h,k}^{\star}
 &=e_k^{\mathsf T}\Delta
 -\widehat\omega_k^{\mathsf T}
 (g^{\star}+H_h\Delta+r_{n,h})\\
 &=-\omega_{k,h}^{\mathsf T}g^{\star}+R_{1k}+R_{2k}+R_{3k},
\end{align*}
where
\begin{align*}
 R_{1k}&=-(\widehat\omega_k-\omega_{k,h})^{\mathsf T}g^{\star},\\
 R_{2k}&=(e_k-H_h\widehat\omega_k)^{\mathsf T}\Delta,\\
 R_{3k}&=-\widehat\omega_k^{\mathsf T}r_{n,h}.
\end{align*}
Theorem~\ref{thm:supp-estimation} and Theorem~\ref{thm:supp-precision} imply
\[
 |R_{1k}|=O_p(s_{\omega,k}\mu_n\lambda_n).
\]
Moreover,
\[
 \norm{e_k-H_h\widehat\omega_k}_{\infty}
 \leq\norm{e_k-\widehat H\widehat\omega_k}_{\infty}
 +\norm{\widehat H-H_h}_{\max}\norm{\widehat\omega_k}_1
 =O_p(\mu_n),
\]
so $|R_{2k}|=O_p(\mu_ns\lambda_n)$. Finally,
$|R_{3k}|=O_p(\mathfrak t_{n,h})$ because
$\norm{\widehat\omega_k}_1=O_p(1)$.
The scaling in Assumption~\ref{ass:supp-expansion} makes each remainder
$o_p(n^{-1/2})$. Since
$g^{\star}=-n^{-1}\sum_iS_{i,h}(\beta_h^{\star})$, the result follows.
\end{proof}

\subsection{Normal approximation and studentisation}\label{supp:normality-proof}

\begin{theorem}[Normality and variance consistency]\label{thm:supp-normality}
Under Assumptions~\ref{ass:supp-estimation}--\ref{ass:supp-variance},
\[
 \frac{\sqrt n(\widetilde\beta_k-\beta_{h,k}^{\star})}{\sigma_k}
 \rightsquigarrow N(0,1),
 \qquad
 \widehat\sigma_k^2\overset{p}{\longrightarrow}\sigma_k^2.
\]
If additionally $\sqrt n h^2\to0$, the statistic may be centred at
$\beta_k^{\star}$.
\end{theorem}

\begin{proof}
The Lyapunov ratio for the independent triangular array
$\{\xi_{ik,h}:i=1,\ldots,n\}$ is bounded by
\[
 \frac{\sum_i\E|\xi_{ik,h}|^{2+\nu}}
 \{\sum_i\Var(\xi_{ik,h})\}^{1+\nu/2}
 \leq C n^{-\nu/2}\to0.
\]
The cluster central limit theorem and Theorem~\ref{thm:supp-bahadur} therefore give the first assertion for $\beta_{h,k}^{\star}$.
Theorem~\ref{thm:supp-target} and $\sqrt n h^2\to0$ permit replacement by $\beta_k^{\star}$.

For variance consistency, write
$\widehat\xi_{ik}=\widehat\omega_k^{\mathsf T}(\widehat S_i-\overline S)$ and
$\xi_{ik,h}=\omega_{k,h}^{\mathsf T}S_{i,h}^{\star}$. Before empirical centring,
\begin{align*}
 \widehat\omega_k^{\mathsf T}\widehat S_i-\xi_{ik,h}
 &=(\widehat\omega_k-\omega_{k,h})^{\mathsf T}S_{i,h}^{\star}
 +\widehat\omega_k^{\mathsf T}(\widehat S_i-S_{i,h}^{\star}).
\end{align*}
The average square of the first term is
$O_p(s_{\omega,k}^2\mu_n^2\log p)=o_p(1)$ under coordinatewise cluster-score moment bounds; the average square of the second is $o_p(1)$ by \eqref{eq:supp-score-plugin}. Empirical centring changes the average squared difference by $o_p(1)$. Consequently,
$n^{-1}\sum_i(\widehat\xi_{ik}-\xi_{ik,h})^2=o_p(1)$.
The cluster law of large numbers gives
$n^{-1}\sum_i\xi_{ik,h}^2\to_p\sigma_k^2$, proving variance consistency.
\end{proof}

\subsection{A feasible asymptotic regime}\label{supp:feasible}

\begin{proposition}[Non-empty scaling region]\label{prop:supp-feasible}
Let $s$ and $s_{\omega,k}$ be bounded, $h=n^{-1/3}$,
$\lambda_n=\sqrt{\log(p)/n}$,
\[
 \eta_{n,h}=O_p(\lambda_n/h),\qquad
 \mathfrak t_{n,h}=O_p(\lambda_n^2/h),\qquad
 \log p=o(n^{1/6}),
\]
and
$\delta_{n,h}=\sqrt{\log(p)/(nh)}+h^2+\eta_{n,h}$,
$\mu_n\asymp\delta_{n,h}$. Then all scaling conditions in
Assumptions~\ref{ass:supp-precision}--\ref{ass:supp-variance}, except the model-specific score plug-in statement \eqref{eq:supp-score-plugin}, are mutually compatible.
\end{proposition}

\begin{proof}
Under $h=n^{-1/3}$,
\[
 \sqrt{\frac{\log p}{nh}}=n^{-1/3}\sqrt{\log p},
 \quad h^2=n^{-2/3},
 \quad \frac{\lambda_n}{h}=n^{-1/6}\sqrt{\log p}.
\]
Thus $\mu_n=O_p(n^{-1/6}\sqrt{\log p})$. It follows that
\[
 \sqrt n\,\mu_n\lambda_n=O_p(n^{-1/6}\log p),
 \quad
 \mu_n\sqrt{\log p}=O_p(n^{-1/6}\log p),
\]
and
\[
 \sqrt n\,\mathfrak t_{n,h}
 =O_p(n^{-1/6}\log p).
\]
All tend to zero when $\log p=o(n^{1/6})$. In addition,
$\sqrt n h^2=n^{-1/6}\to0$ and
$\mu_n^2\log p=O_p(n^{-1/3}(\log p)^2)=o_p(1)$.
\end{proof}

